\documentclass[11pt,twoside,a4paper]{article}
%NOTE: This report format is 

\newcommand{\reporttitle}{Group 23: {Your title}}
\newcommand{\reportauthorOne}{Pang Jia De}
\newcommand{\cidOne}{2402226, 3070684P}
\newcommand{\reportauthorTwo}{Daryl Chong Han Kin}
\newcommand{\cidTwo}{2400619, 3070590C}
\newcommand{\reportauthorThree}{Kaiser Yee}
\newcommand{\cidThree}{2400821, 3070609Y}
\newcommand{\reporttype}{Coursework}
\bibliographystyle{plain}

% include files that load packages and define macros
\input{includes} % various packages needed for maths etc.
\input{notation} % short-hand notation and macros


%===================================================================

\begin{document}
% front page
\input{titlepage}


%===================================================================
% TABLE OF CONTENT
%===================================================================
\tableofcontents
\newpage

%=====================================================================
% 1. OVERALL PROJECT DESCRIPTION  (3–6 pages, part of 20% criterion)
%=====================================================================
\section{Introduction}

\subsection{Background and Motivation}
Earth-observation satellites, aerial surveys, and high-resolution basemaps now produce far more imagery than analysts can review by hand \cite{Zhu2017}. Getting useful information out of that imagery like where infrastructure is located, how an area is changing, matters for urban planning, environmental monitoring, infrastructure management, and disaster response, among other applications. Scene classification, labelling an image patch with a land-use category, is the basic task that most of this depends on.

Until the early 2010s this problem was tackled with handcrafted features: colour histograms, texture descriptors, SIFT-based bag-of-visual-words. These worked reasonably well for coarse distinctions but stalled because they describe low-level appearance rather than the structure of a scene. Deep convolutional networks changed that by learning features directly from data instead of relying on hand-designed ones. Since AlexNet's ImageNet result \cite{Krizhevsky2012}, they have taken over image classification generally, and the remote-sensing literature shows the same pattern: deep models are now standard on aerial scene classification benchmarks \cite{Cheng2020}.

The dataset used in this project makes the task harder in two specific ways. First, the four classes: \texttt{coastal\_mansion}, \quad \texttt{dense\_residential}, \quad \texttt{nursing\_home}, \quad \texttt{sparse\_residential}, all seen from above look visually similar to each other. What separates them is density and spatial arrangement rather than distinct objects or colours, so a model has to pick up on layout, not just presence of things. Second, there are only 700 training images per class, which is little enough that a network trained from scratch would likely overfit. This is why data augmentation and transfer learning from pre-trained models are central to our approach rather than optional extras.

We also cared about two things beyond accuracy. One is fair comparison: architecture papers are often hard to compare because they differ in input resolution, augmentation, training schedule, or evaluation setup, so we fixed all of these across our experiments and varied only the architecture. The other is deployment; a model is only useful if it can actually be run somewhere. So we packaged the best-performing model as a containerized inference service rather than leaving it as a notebook result.

\subsection{Problem Statement}
Each team in this cohort is assigned a partition of a 38-category aerial-imagery benchmark. The class names, the 800 images per class, and the native $256 \times 256$ pixel resolution identify the source as \textbf{PatternNet}, a high-resolution dataset collected mainly over US cities from Google Earth and Google Map API imagery. Our partition (\emph{set 23}) contains four categories:
\begin{center}
\texttt{coastal\_mansion}, \quad \texttt{dense\_residential}, \quad
\texttt{nursing\_home}, \quad \texttt{sparse\_residential},
\end{center}
each with 700 training images and 100 held-out validation images. The held-out
images never appear in training and are used only for final evaluation.

The task is standard four-way image classification: train a model that takes a
$256 \times 256$ RGB aerial image and outputs a probability for each of the
four classes, and evaluate it on the held-out set using accuracy, per-class
and macro-averaged precision, recall, and $F_1$-score, and the confusion
matrix.

What makes this partition hard is that all four categories are residential or
institutional building scenes. A nursing home is essentially a large
residential complex; dense and sparse residential areas differ in lot size and
building spacing rather than in kind; the most reliable cue for a coastal
mansion is the shoreline next to it, not the building itself. So the classes
look alike (same mix of buildings, roads, and vegetation) while images within
a class vary a lot (roof colour, vegetation cover, season, layout).

\begin{figure}[htbp]
    \centering

    % Column labels
    \begin{minipage}[b]{0.23\textwidth}
        \centering \texttt{coastal\_mansion}
    \end{minipage}\hfill
    \begin{minipage}[b]{0.23\textwidth}
        \centering \texttt{dense\_residential}
    \end{minipage}\hfill
    \begin{minipage}[b]{0.23\textwidth}
        \centering \texttt{nursing\_home}
    \end{minipage}\hfill
    \begin{minipage}[b]{0.23\textwidth}
        \centering \texttt{sparse\_residential}
    \end{minipage}

    \vspace{0.3em}

    % Row 1
    \begin{minipage}[b]{0.23\textwidth}
        \includegraphics[width=\linewidth]{docs/figures/coastalmansion001.jpg}
    \end{minipage}\hfill
    \begin{minipage}[b]{0.23\textwidth}
        \includegraphics[width=\linewidth]{docs/figures/denseresidential001.jpg}
    \end{minipage}\hfill
    \begin{minipage}[b]{0.23\textwidth}
        \includegraphics[width=\linewidth]{docs/figures/nursinghome001.jpg}
    \end{minipage}\hfill
    \begin{minipage}[b]{0.23\textwidth}
        \includegraphics[width=\linewidth]{docs/figures/sparseresidential001.jpg}
    \end{minipage}

    \vspace{0.4em}

    % Row 2
    \begin{minipage}[b]{0.23\textwidth}
        \includegraphics[width=\linewidth]{docs/figures/coastalmansion002.jpg}
    \end{minipage}\hfill
    \begin{minipage}[b]{0.23\textwidth}
        \includegraphics[width=\linewidth]{docs/figures/denseresidential002.jpg}
    \end{minipage}\hfill
    \begin{minipage}[b]{0.23\textwidth}
        \includegraphics[width=\linewidth]{docs/figures/nursinghome002.jpg}
    \end{minipage}\hfill
    \begin{minipage}[b]{0.23\textwidth}
        \includegraphics[width=\linewidth]{docs/figures/sparseresidential002.jpg}
    \end{minipage}

    \caption{Representative samples from each class in our partition (set 23) of PatternNet. All four categories share buildings, roads, and vegetation as their basic visual vocabulary; the distinguishing cues are density, spacing, and, for coastal mansions, shoreline adjacency rather than the built structures themselves.}
    \label{fig:class-samples}
\end{figure}

The project specification adds three practical requirements: (i) the team must
investigate at least as many deep learning approaches as it has members
(five); (ii) every open-source library used must be registered and justified,
and our own contributions documented; and (iii) the best model must be
deployed in a container that accepts an aerial image and returns a predicted
label with confidence scores.

\subsection{Project Objectives}
We set ourselves six objectives:

\begin{enumerate}
  \item \textbf{O1: Understand the data.} Explore class balance, image
  statistics, and visual overlap between categories, and use the findings to
  justify pre-processing and augmentation choices.
  \item \textbf{O2: Evaluate without leakage.} Use a shared data pipeline
  with a fixed internal tuning split for development and checkpoint selection,
  and touch the 400 held-out images only once, for final evaluation.
  \item \textbf{O3: Investigate five approaches.} One per team member: a
  from-scratch CNN baseline, three pretrained CNN families (ResNet-50,
  EfficientNet-B0, MobileNetV2), and a vision transformer.
  \item \textbf{O4: Compare fairly.} Train every model with the same seed,
  input size, augmentation, batch size, epoch budget, and early-stopping
  patience, so that differences in results come from the architectures.
  \item \textbf{O5: Maximise held-out performance.} Get the best accuracy
  and macro-$F_1$ we can, and use confusion matrices to understand the
  remaining errors.
  \item \textbf{O6: Deploy the best model.} Package it in a Docker
  container with a web interface that returns the predicted category and
  confidence scores, with automated contract tests covering the shared data
  pipeline and training entrypoint.
\end{enumerate}

\subsection{Review of Existing Approaches}
\label{sec:review}

\subsubsection{From Handcrafted Descriptors to Learned Features}

Early remote-sensing scene classification relied on handcrafted descriptors ---
colour histograms, texture measures, and SIFT-based bag-of-visual-words
representations. These encode low-level appearance and therefore struggle
precisely where our partition is hardest: distinguishing categories that share
the same visual vocabulary and differ only in spatial arrangement. The shift to
learned representations after AlexNet \cite{Krizhevsky2012} resolved this by
inferring features directly from data, and the remote-sensing literature has
followed the same trajectory \cite{Zhu2017, Cheng2020}.

The magnitude of the improvement is well documented. Kwenda et al.
\cite{Kwenda2024} isolate it cleanly on aerial forest imagery: classical
classifiers applied to raw pixels reach 0.646 accuracy, whereas the same
classifiers applied to features extracted by a pretrained deep network reach
0.940. Zhou et al. \cite{Zhou2018} report the same ordering on the benchmark
used in this project, where CNN-derived features outperform handcrafted
descriptors for retrieval. Both results motivate our decision to treat deep
feature extraction, rather than feature engineering, as the foundation of every
approach we investigate.

\subsubsection{Benchmark Datasets}

Progress in this field has been driven by successive benchmarks --- UC Merced,
WHU-RS19, RSSCN7, AID, and NWPU-RESISC45 --- each larger or more finely
categorised than the last. Our data comes from \textbf{PatternNet}
\cite{Zhou2018}, a 38-class collection of 800 images per class at
$256 \times 256$ pixels, gathered over US cities via Google Earth imagery and
the Google Map API.

One design property of PatternNet directly shapes this project. The dataset was
constructed for \emph{image retrieval}, a task that rewards fine-grained
category distinctions, so it deliberately separates classes that a
coarser-grained benchmark would merge. Our partition inherits four such
categories --- \texttt{coastal\_mansion}, \texttt{dense\_residential},
\texttt{nursing\_home}, and \texttt{sparse\_residential} --- all of which are
residential or institutional building scenes. The difficulty we report in
Section~\ref{sec:results} is therefore intrinsic to the benchmark's design
rather than incidental to our partition.

\subsubsection{Convolutional Architectures}

The systematic review by Shoaib et al. \cite{Shoaib2025}, covering 124 studies
of deep learning for aerial imagery, identifies convolutional networks as the
dominant family for classification, valued for accuracy and computational
efficiency. Three developments in that family inform our choices.

\textbf{Residual learning.} ResNet \cite{He2016} introduced skip connections so
that each block learns a residual adjustment to its input, allowing gradients
to propagate through very deep stacks. ResNet-50 remains a standard baseline
with well-understood transfer behaviour, which is why we include it.

\textbf{Compound scaling.} EfficientNet \cite{Tan2019} observes that network
depth, width, and input resolution should be scaled together under a single
coefficient rather than independently, yielding markedly better accuracy per
parameter. Shoaib et al. \cite{Shoaib2025} single out this efficiency as
important for high-resolution aerial data, where computational cost scales
quickly with image size.

\textbf{Efficiency-oriented design.} MobileNetV2 \cite{Sandler2018} factorises
standard convolutions into depthwise and pointwise operations and adds inverted
residual blocks with linear bottlenecks, reducing parameter count by an order of
magnitude at modest accuracy cost. It represents the deployment-constrained end
of the design space, and the development of such lightweight models is
identified as a priority direction in \cite{Shoaib2025}.

\subsubsection{Transfer Learning under Limited Supervision}

With 700 training images per class, training a large network from random
initialisation risks overfitting. The standard remedy is transfer learning:
reusing representations learned on a large source corpus such as ImageNet.
Shoaib et al. \cite{Shoaib2025} state the case directly --- pretrained models
``provide robust feature representations that generalize well to aerial
imagery,'' and this ``is especially important for scenarios with limited
labeled data.''

The practical question is how much of the network to adapt. Kwenda et al.
\cite{Kwenda2024} take the conservative extreme, freezing VGG16 and ResNet-50
entirely and training only a downstream classifier, and still obtain strong
results. We adopt the same strategy for our three pretrained models: the
backbone is frozen and only a small classification head is fitted. This
suppresses overfitting on a small dataset and, equally important for our
purpose, holds the amount of adaptation constant across the transfer-learning
approaches so that the comparison reflects architecture rather than fine-tuning
effort.

\subsubsection{Attention and Vision Transformers}

Vision Transformers \cite{Dosovitskiy2021} discard convolution entirely,
partitioning an image into fixed patches and processing the resulting sequence
with self-attention. Every patch can attend to every other from the first
layer, capturing global spatial layout that a convolutional stack only reaches
after many layers --- a property well matched to aerial scenes, where the
discriminating cue is the arrangement of elements across the whole image.

This capability carries a known cost. Shoaib et al. \cite{Shoaib2025} caution
that transformers ``typically require large amounts of training data and
computational resources, which can be a limitation in domains like aerial image
analysis with limited annotated datasets.'' Whether that limitation binds at our
dataset scale is one of the questions this study answers empirically
(Section~\ref{sec:comparative}).

\subsubsection{Positioning of This Work}

The literature establishes that deep features outperform handcrafted ones, that
transfer learning is the standard response to limited data, and that
architectures differ in their accuracy-to-cost trade-offs. What it does not
readily provide is a \emph{controlled} comparison: published results differ in
input resolution, augmentation, optimiser, training schedule, and evaluation
protocol, so architectural claims are difficult to isolate. This project
addresses that gap on one dataset by fixing every such variable across five
architectures and varying only the model, as detailed in
Section~\ref{sec:splits} and Section~\ref{sec:augmentation}.

\subsubsection{Open-Source Libraries Used}

Table~\ref{tab:libraries} registers every third-party component used in this
project. All are open-source and used under their respective licences; no
pretrained checkpoint or code was used outside its intended terms.

\begin{table}[H]
\centering
\caption{Open-source libraries and tools, with the role each plays.}
\label{tab:libraries}
\small
\begin{tabular}{lll}
\toprule
Library & Version & Role in this project \\
\midrule
TensorFlow / Keras & $\geq$ 2.15 & Model definition, training, augmentation layers, checkpointing \\
keras-hub & $\geq$ 0.3.0 & Pretrained Vision Transformer backbone \\
scikit-learn & $\geq$ 1.3 & Precision, recall, $F_1$, and confusion-matrix computation \\
NumPy & $\geq$ 1.24 & Array operations and metric aggregation \\
pandas & $\geq$ 2.0 & Tabular summaries during exploratory analysis \\
Matplotlib & $\geq$ 3.7 & Figures: sample grids, training curves, confusion matrices \\
Pillow & $\geq$ 10.0 & Image loading and inspection for EDA \\
pytest & $\geq$ 8.0 & Automated contract tests for the data and training pipelines \\
Streamlit & 1.58.0 & Web interface for the deployed classifier \\
Docker & --- & Containerisation of the deployed model (\texttt{python:3.12-slim}) \\
\bottomrule
\end{tabular}
\end{table}

Pretrained ImageNet weights for ResNet-50, EfficientNet-B0, and MobileNetV2 are
obtained through the Keras applications module; the Vision Transformer backbone
is obtained through keras-hub. All model definitions, the shared data pipeline,
the augmentation policy, the training and evaluation entrypoints, the automated
tests, and the deployment application were written by the team; our specific
contributions and modifications are documented in
Section~\ref{sec:contributions}.

%=====================================================================
% 2. MACHINE LEARNING SOLUTIONS  (15–18 pages, 60% criterion)
%=====================================================================
\section{Exploratory Data Analysis}
% Source: eda.ipynb. Every figure and statistic below is computed from the
% TRAINING split only (plus validation *counts*); no modelling decision is
% informed by the held-out images.

Exploratory analysis was carried out before any model was built, and
deliberately restricted to the training split. Only the \emph{counts} of the
held-out directory were inspected; no held-out image was viewed or measured,
so none of the design decisions that follow could have been informed by the
data used for final evaluation.

\subsection{Class Balance}

The dataset is exactly balanced. Each of the four categories contributes 700
training images and 100 held-out images, giving 2{,}800 and 400 respectively
(Table~\ref{tab:class-balance}).

\begin{table}[H]
\centering
\caption{Class distribution across the provided splits.}
\label{tab:class-balance}
\small
\begin{tabular}{lrrr}
\toprule
Class & Train & Held-out & Total \\
\midrule
\texttt{coastal\_mansion}    & 700 & 100 & 800 \\
\texttt{dense\_residential}  & 700 & 100 & 800 \\
\texttt{nursing\_home}       & 700 & 100 & 800 \\
\texttt{sparse\_residential} & 700 & 100 & 800 \\
\midrule
Total & 2{,}800 & 400 & 3{,}200 \\
\bottomrule
\end{tabular}
\end{table}

Perfect balance has two consequences for the rest of this report. First, plain
accuracy is not misleading here, because a majority-class predictor is capped
at 25\%; we nevertheless report macro-$F_1$ alongside it, since macro-averaging
exposes a collapse on any single class that accuracy alone can hide. Second,
no class weighting or resampling is required, which removes a confounding
variable from the five-model comparison.

\subsection{Image Properties}

Sampling 100 images per class (400 in total) shows the collection to be
perfectly uniform: every image is $256 \times 256$ pixels with zero variance,
an aspect ratio of exactly $1.0$, and all 400 sampled images are in RGB mode.
No greyscale, CMYK, or alpha-channel images are present, and no corrupt files
were encountered.

This uniformity motivated a concrete team decision. The pipeline was
originally configured at the conventional ImageNet resolution of
$224 \times 224$, but because the source images are natively $256 \times 256$,
that setting discarded roughly 23\% of the available pixels before training
began. The shared configuration was therefore standardised at
\texttt{IMAGE\_SIZE = 256} so that models operate on the data at its native
resolution and no downsampling artefacts are introduced. The trade-off --- the
ImageNet-pretrained backbones no longer receive their native input size --- is
revisited in Section~\ref{sec:comparative}.

\subsection{Colour Statistics}

Mean per-channel intensities were computed over 60 images per class
(Table~\ref{tab:colour-signature}). The striking result is how little colour
separates the categories: all twelve channel means fall within a narrow band
of roughly 60--100 intensity units.

\begin{table}[H]
\centering
\caption{Mean per-channel intensity by class (60 training images per class).}
\label{tab:colour-signature}
\small
\begin{tabular}{lrrr}
\toprule
Class & R & G & B \\
\midrule
\texttt{coastal\_mansion}    & 74.2 & 71.6 & 63.0 \\
\texttt{dense\_residential}  & 89.7 & 86.3 & 77.3 \\
\texttt{nursing\_home}       & 96.9 & 100.6 & 88.0 \\
\texttt{sparse\_residential} & 70.4 & 74.6 & 60.7 \\
\bottomrule
\end{tabular}
\end{table}

\texttt{coastal\_mansion} and \texttt{sparse\_residential} are especially
close, differing by under 4 intensity units in every channel --- both are
low-density scenes dominated by vegetation and large plots.
\texttt{nursing\_home} is the brightest class, consistent with large
institutional roofs. The practical implication is that a classifier cannot
succeed on colour alone: the discriminative signal lies in \emph{texture and
spatial structure} --- building density, roof arrangement, road layout, and the
regularity of the built pattern. This prediction is borne out by the results in
Section~\ref{sec:results}, where the residual errors are concentrated in
exactly the density-based distinctions rather than colour-based ones.

\subsection{Visual Confusability}

Inspecting sample grids per class confirms that this is a genuinely
fine-grained problem rather than a coarse one. All four categories are
suburban aerial scenes photographed at similar altitude, sharing the same
vocabulary of roofs, roads, and vegetation; they differ only in the
\emph{arrangement} and \emph{density} of those elements. Two pairs are
particularly difficult:

\begin{itemize}
  \item \texttt{dense\_residential} versus \texttt{sparse\_residential} --- an
    identical scene type separated only by a continuous, unlabelled density
    threshold, with no discrete visual boundary between them.
  \item \texttt{nursing\_home} versus \texttt{dense\_residential} --- both show
    clustered multi-unit buildings; the distinction rests on building
    footprint and layout regularity rather than on scene content.
\end{itemize}

These predictions were recorded before training and are directly confirmed by
the confusion matrices in Section~\ref{sec:results}.

\subsection{Data Quality: Duplicate Images}

One integrity problem was found. Two training files are byte-identical to each
other yet assigned to conflicting classes:

\begin{center}
\texttt{coastal\_mansion/coastalmansion029.jpg} \quad and \quad
\texttt{sparse\_residential/sparseresidential021.jpg}
\end{center}

The same image therefore carries two mutually exclusive labels. Left in place
it supplies contradictory gradients and, worse, could place identical pixels on
both sides of the train/tuning split --- a small but genuine leakage path.
Both files are excluded logically by the shared loader through
\texttt{EXCLUDED\_TRAIN\_FILES} in \texttt{shared/config.py}; the source data on
disk is never modified, so the exclusion is reproducible and auditable by every
team member. The usable training pool is consequently 2{,}798 images. Notably,
the affected pair belongs to the two classes that colour statistics already
identified as most similar.

\subsection{Implications for Modelling}

Four conclusions carried forward into the design:

\begin{enumerate}
  \item \textbf{Accuracy is a fair headline metric} (balanced classes), but
    macro-$F_1$ is reported alongside it to expose per-class collapse.
  \item \textbf{Inputs are standardised at $256 \times 256$}, the native
    resolution, avoiding avoidable information loss.
  \item \textbf{Texture and structure, not colour, carry the signal}, so the
    augmentation policy must preserve fine-grained texture --- photometric
    jitter is kept deliberately mild for this reason
    (Section~\ref{sec:augmentation}).
  \item \textbf{Aerial scenes have no canonical orientation.} Unlike natural
    photographs, a top-down scene is equally valid flipped or rotated, so full
    horizontal and vertical flips and arbitrary rotations are
    \emph{label-preserving} here. This observation directly justifies the
    shared augmentation policy.
\end{enumerate}

\section{Data Preparation and Pre-Processing}
\subsection{Data Splits and Leakage Control}
\label{sec:splits}

A recurring flaw in student-scale deep learning experiments is that the same
data guides both model selection and final reporting, which inflates the
headline number without anyone intending to cheat. We designed the pipeline
specifically to make that mistake impossible, using a three-way split
implemented once in \texttt{shared/data.py} and shared by all five models.

\paragraph{Three-way split.}
The provided \texttt{val 23} directory is treated as a \emph{sealed test set}.
All development data is therefore carved out of \texttt{set 23} alone: a
deterministic stratified split assigns 85\% of it to model fitting and 15\% to
an internal \emph{tuning} set (Table~\ref{tab:splits}).

\begin{table}[H]
\centering
\caption{Three-way data split. Only the training subset updates weights; the
tuning subset guides checkpoint selection; the held-out set is untouched until
final evaluation.}
\label{tab:splits}
\small
\begin{tabular}{llrl}
\toprule
Subset & Source & Images & Role \\
\midrule
Training & \texttt{set 23} & 2{,}378 & gradient updates \\
Tuning   & \texttt{set 23} & 420 & early stopping, checkpoint selection \\
Held-out & \texttt{val 23} & 400 & final evaluation only \\
\bottomrule
\end{tabular}
\end{table}

The split is \emph{stratified}: it is drawn independently within each class, so
the tuning subset contains exactly 105 images per class and the training
subset 594--595 per class. Class balance is thus preserved on both sides rather
than left to chance, which matters because per-class $F_1$ on a 420-image
subset is sensitive to even small imbalances.

\paragraph{Determinism.}
Reproducibility is enforced by seeding the shuffle for class $k$ with
$\texttt{SEED} + k$ (\texttt{SEED} $=42$) rather than by a global random state.
Consequently the partition depends only on the file listing --- not on
execution order, machine, or how many random draws happened earlier in the
session. Every team member, on every machine, trains on byte-identical
subsets. The exact counts are recorded in each run's metadata via
\texttt{get\_split\_counts()}, so any run can be audited after the fact.

\paragraph{Why the held-out set must stay sealed.}
Both \texttt{ModelCheckpoint} and \texttt{EarlyStopping} select the epoch that
maximises \emph{tuning} accuracy. That makes tuning metrics optimistically
biased by construction: they measure the best of roughly 30 attempts, not the
expected performance on new data. The effect is not hypothetical --- for our
custom CNN the tuning macro-$F_1$ of 0.908 overstates the held-out 0.845 by
more than six points. Every number reported in
Section~\ref{sec:results} is therefore held-out, and the tuning metrics are
retained only as a training diagnostic.

Mechanically, the held-out loader is a separate function that reads
\texttt{val 23} and is invoked only behind an explicit
\texttt{--evaluate-held-out} flag, after training has finished. It is never
passed to \texttt{model.fit} and never influences a callback. Combined with the
duplicate-image exclusion described above --- which prevents identical pixels
from appearing on both sides of the split --- no image can reach evaluation
having been seen during training.

\paragraph{Deviation from the literature protocol.}
Our effective proportions are approximately 74/13/13 rather than the 70/15/15
used by the reference architecture we follow, because the test partition is
fixed at 400 images by the provided data rather than chosen by us. Since the
identical split is applied to all five models, this affects absolute values
only, not the comparison between approaches.
\subsection{Input Resolution}

All models receive $256 \times 256$ RGB inputs. As established in the
exploratory analysis, this is the \emph{native} resolution of every image in
the collection, so no resizing artefact is introduced and no information is
discarded. The pipeline was initially configured at the conventional ImageNet
resolution of $224 \times 224$, which would have thrown away 23.4\% of the
available pixels before training began; standardising on 256 removes that loss.

The trade-off is that the ImageNet-pretrained backbones no longer operate at
the resolution they were trained on. We accepted this deliberately: the
alternative penalises the from-scratch model (which has no such prior) in order
to flatter the pretrained ones, and comparability requires that all five
approaches see identical inputs. Resizing uses bilinear interpolation with
antialiasing enabled, which matters when a model internally rescales.

\subsection{Augmentation}
\label{sec:augmentation}

With only 2{,}378 training images against models of 2.3M--85.8M parameters,
augmentation is not an optional refinement but the main defence against
overfitting. The policy is defined once in \texttt{shared/augment.py} and
embedded as the first stage of every model, so it is applied identically across
all five approaches and is automatically inactive at inference time.

\begin{table}[H]
\centering
\caption{Shared training-time augmentation policy. All layers are seeded for
reproducibility and operate on raw $[0,255]$ pixels.}
\label{tab:augmentation}
\small
\begin{tabular}{lll}
\toprule
Transform & Setting & Justification \\
\midrule
Random flip & horizontal + vertical & no canonical orientation \\
Random rotation & $\pm 180^\circ$, reflect fill & no canonical orientation \\
Random zoom & $\pm 20\%$ & capture altitude / scale variation \\
Random contrast & $\pm 10\%$ & sensor and lighting variation \\
Random brightness & $\pm 10\%$ & sensor and lighting variation \\
\bottomrule
\end{tabular}
\end{table}

\paragraph{Geometric transforms are label-preserving here.}
For natural photographs a vertical flip produces an implausible image, and
practitioners routinely restrict augmentation to horizontal flips and small
rotations for that reason. Top-down aerial imagery has no such constraint: a
residential block photographed from above is equally valid rotated by any
angle or mirrored in either axis, and the class label is unchanged. Full
flips and arbitrary rotations are therefore \emph{free} additional training
signal rather than label noise, and we exploit them aggressively. Rotation
uses \texttt{reflect} fill so that rotated corners contain plausible texture
instead of black wedges, which would otherwise be a spurious cue correlated
with augmentation rather than with class.

\paragraph{Photometric transforms are kept deliberately mild.}
Contrast and brightness are restricted to $\pm 10\%$ --- an order of magnitude
weaker than the geometric transforms. This asymmetry is a direct consequence
of the exploratory finding that the four classes are \emph{not} separable by
colour: with all twelve channel means confined to a narrow band, the
discriminative signal lies in fine-grained texture and spatial structure.
Aggressive photometric jitter would degrade precisely the cues the models
depend on. This was confirmed empirically --- an early configuration using a
$\pm 50\%$ factor drove images to near-black or near-white and destroyed the
texture detail entirely, so the factor was reduced.

\paragraph{Ordering and value range.}
The augmentation layers operate on raw $[0,255]$ pixels and return values in
the same range, so any model-specific normalisation (a \texttt{Rescaling}
layer for the custom CNN, or the backbone's own \texttt{preprocess\_input} for
the transfer models) is applied \emph{after} augmentation. Placing rescaling
first is a subtle but real bug: the brightness layer is configured for a
$[0,255]$ value range, and feeding it $[0,1]$ inputs saturates the image to
black. Because augmentation lives inside the model graph rather than in the
data pipeline, this ordering is fixed at construction time and cannot be got
wrong at inference.

\paragraph{Augmentation is training-only.}
The Keras preprocessing layers are inactive outside training, and the tuning
and held-out datasets are never passed through the augmentation pipeline. All
reported metrics therefore reflect performance on clean, untouched images.

\paragraph{Status as a controlled variable.}
The policy is fixed and shared rather than tuned per model. Individual
architectures might well prefer slightly different settings, but allowing each
member to tune augmentation independently would confound the comparison: a
difference in final accuracy could then reflect a better augmentation search
rather than a better architecture. Holding augmentation constant --- along with
the data, split, input resolution, and evaluation procedure --- ensures that
the leaderboard in Section~\ref{sec:results} isolates the effect of
architecture, which is the stated objective of this study.

\section{Deep Learning Approaches}
% One subsection per approach (>= number of team members = 5).
% For EACH: architecture diagram/description, design decisions, transfer
% learning strategy (frozen backbone vs fine-tuned; note trainable vs total
% parameter counts), and what was customised — the spec penalises
% "download and run" without meaningful modification.
\subsection{Custom CNN (from scratch)}
\label{sec:customcnn}
% 26.5M params, all trainable. Serves as the no-pretraining baseline.

The custom CNN is the only approach in this study that uses no pretrained
weights. Every one of its 26,512,212 parameters is learned from our 2,378
training images alone. Its role is therefore twofold: it is a competitive
architecture in its own right, and it is the experimental control that makes
the contribution of ImageNet pretraining measurable rather than assumed --- a
comparison the other four approaches cannot provide on their own.

\subsubsection{Source Architecture}

Rather than assemble an arbitrary stack of layers, we reproduce the
\emph{Pyramidal Net} architecture proposed by Abdelmonem et al.
\cite{Abdelmonem2026}, which was designed specifically for remote-sensing scene
classification and reports 94.28\% accuracy on a ten-class NWPU-RESISC45 subset
and 93\% on UC~Merced. The choice is well matched to our setting: the source
study uses 700 images per class --- the same per-class budget as our partition
--- and operates natively at $256 \times 256$, so its design assumptions
transfer to our data without reinterpretation.

The name describes the design principle: spatial dimensions contract while
feature abstraction increases, so the representation narrows toward the
classifier like a pyramid. Table~\ref{tab:customcnn-arch} gives the structure as
implemented.

\begin{table}[H]
\centering
\caption{Custom CNN (Pyramidal Net) architecture, following \cite{Abdelmonem2026}.}
\label{tab:customcnn-arch}
\small
\begin{tabular}{lll}
\toprule
Stage & Configuration & Output \\
\midrule
Input & $256 \times 256 \times 3$, raw $[0,255]$ RGB & $256^2 \times 3$ \\
Augmentation & shared team policy, training only & $256^2 \times 3$ \\
Rescaling & $\times \, 1/255$ & $256^2 \times 3$ \\
Conv 1 & 256 filters, $3\times3$, ReLU, valid & $254^2 \times 256$ \\
MaxPool 1 + BatchNorm & $2\times2$, valid & $127^2 \times 256$ \\
Conv 2 & 128 filters, $3\times3$, ReLU, valid & $125^2 \times 128$ \\
MaxPool 2 + BatchNorm & $2\times2$, valid & $62^2 \times 128$ \\
Conv 3 & 128 filters, $3\times3$, ReLU, valid & $60^2 \times 128$ \\
Conv 4 & 128 filters, $3\times3$, ReLU, valid & $58^2 \times 128$ \\
Conv 5 & 64 filters, $3\times3$, ReLU, valid & $56^2 \times 64$ \\
MaxPool 3 & $2\times2$, valid & $28^2 \times 64$ \\
Flatten & --- & 50{,}176 \\
Dense block & $512 \rightarrow 256 \rightarrow 64 \rightarrow 16$, ReLU & 16 \\
Output & Dense, softmax & 4 \\
\midrule
\multicolumn{2}{l}{Parameters} & 26{,}512{,}212 (all trainable) \\
\bottomrule
\end{tabular}
\end{table}

The flattened vector of 50,176 features matches the source paper exactly,
confirming a faithful reproduction of the convolutional stack. All five
convolutional layers use $3\times3$ kernels, unit stride, \texttt{valid}
padding, and ReLU activation, and batch normalisation follows the first two
pooling stages, all as specified in \cite{Abdelmonem2026}.

\subsubsection{Training Configuration}

The optimisation settings follow the source paper's Table~4. The learning-rate
schedule in particular is specific to this architecture and is supplied through
the shared training entrypoint's per-model callback hook, so it applies to this
model alone and does not alter any teammate's run.

\begin{table}[H]
\centering
\caption{Custom CNN training configuration, after \cite{Abdelmonem2026} Table~4.}
\label{tab:customcnn-config}
\small
\begin{tabular}{lll}
\toprule
Hyperparameter & Source paper & This work \\
\midrule
Optimiser & Adam & Adam \\
Initial learning rate & 0.00025 & 0.00025 \\
LR scheduler & ReduceLROnPlateau & ReduceLROnPlateau \\
LR reduction factor & 0.25 & 0.25 \\
Scheduler patience & 4 epochs & 4 epochs \\
Minimum learning rate & $1 \times 10^{-6}$ & $1 \times 10^{-6}$ \\
Weight decay & L2 norm term & L2 $(10^{-4})$ on conv layers \\
Dropout & none & none \\
Kernel / stride / padding & $3\times3$ / 1 / valid & identical \\
Activation & ReLU & ReLU \\
Loss & sparse categorical CE & categorical CE (one-hot) \\
Epochs & 30 & 32 \\
Batch size & 8 & 16 (team standard) \\
Random seed & 32 & 42 (team standard) \\
Split & 70/15/15 & 74/13/13 (fixed test set) \\
Normalisation & resize only & resize + rescale to $[0,1]$ \\
\bottomrule
\end{tabular}
\end{table}

Four deviations are worth justifying explicitly. First, the loss is categorical
cross-entropy over one-hot labels rather than the sparse variant; these are
mathematically identical objectives differing only in label encoding, and the
one-hot form is what the shared data pipeline emits for all five models.
Second, batch size, seed, and split are dictated by the team-wide controlled
variables described in Section~\ref{sec:splits}, without which the five
approaches would not be comparable. Third, the source applies no pixel
normalisation, keeping inputs in $[0,255]$; we insert a rescaling layer after
augmentation because a network trained from random initialisation optimises
more stably on normalised inputs, and the batch-normalisation layers make the
change largely self-correcting in any case. Fourth, the paper specifies weight
decay as ``an L2 norm term'' without a coefficient; we apply $10^{-4}$ to the
convolutional layers, leaving the dense head unregularised.

\subsubsection{Results}

Trained under the shared protocol at batch size 16 for 32 epochs, the model
selected its best checkpoint at epoch 24 and achieved \textbf{84.50\% accuracy}
and \textbf{0.8452 macro-$F_1$} on the held-out set --- the lowest of our five
approaches. The per-class breakdown localises the shortfall precisely:

\begin{center}
\small
\begin{tabular}{lc}
\toprule
Class & $F_1$ \\
\midrule
\texttt{coastal\_mansion} & 0.990 \\
\texttt{nursing\_home} & 0.864 \\
\texttt{sparse\_residential} & 0.800 \\
\texttt{dense\_residential} & \textbf{0.727} \\
\bottomrule
\end{tabular}
\end{center}

Performance is essentially saturated on \texttt{coastal\_mansion}, whose
shoreline provides an unambiguous global cue, and collapses on
\texttt{dense\_residential}. The confusion matrix in
Figure~\ref{fig:customcnn-cm} shows that this is not diffuse error but a single
dominant confusion: \texttt{dense\_residential} is misclassified as
\texttt{sparse\_residential} in a large fraction of cases, while
\texttt{sparse\_residential} itself is recovered almost perfectly. The error is
therefore \emph{asymmetric} --- the model under-predicts density rather than
confusing the two classes symmetrically.

\begin{figure}[htbp]
\centering
\begin{minipage}[t]{0.49\textwidth}
  \centering
  \includegraphics[width=\linewidth]{docs/figures/custom_cnn_tuning_confusion_matrix.png}
  \\[0.2em] {\small (a) Tuning split (420 images)}
\end{minipage}\hfill
\begin{minipage}[t]{0.49\textwidth}
  \centering
  \includegraphics[width=\linewidth]{docs/figures/custom_cnn_confusion_matrix.png}
  \\[0.2em] {\small (b) Held-out set (400 images)}
\end{minipage}
\caption{Custom CNN confusion matrices. (a) The tuning split, on which the
checkpoint was selected. (b) The sealed held-out set. Errors concentrate
almost entirely in the \texttt{dense\_residential} $\rightarrow$
\texttt{sparse\_residential} direction, and that confusion is markedly worse on
the held-out data --- the visual signature of checkpoint-selection bias.}
\label{fig:customcnn-cm}
\end{figure}

This is exactly the failure mode the exploratory analysis predicted: the two
classes share an identical visual vocabulary and are separated only by a
continuous density threshold, so the distinction requires fine-grained texture
discrimination rather than scene-level cues.

\paragraph{Checkpoint-selection bias is class-specific.}
Comparing the two panels of Figure~\ref{fig:customcnn-cm} quantifies the
optimism introduced by selecting the checkpoint on tuning accuracy, and shows
that it is not spread evenly across classes (Table~\ref{tab:customcnn-bias}).

\begin{table}[H]
\centering
\caption{Per-class $F_1$ on the tuning split versus the held-out set.}
\label{tab:customcnn-bias}
\small
\begin{tabular}{lccr}
\toprule
Class & Tuning & Held-out & $\Delta$ \\
\midrule
\texttt{coastal\_mansion}    & 0.950 & 0.990 & $+0.040$ \\
\texttt{nursing\_home}       & 0.871 & 0.864 & $-0.007$ \\
\texttt{sparse\_residential} & 0.893 & 0.800 & $-0.093$ \\
\texttt{dense\_residential}  & 0.917 & 0.727 & $\mathbf{-0.190}$ \\
\midrule
Macro-$F_1$ & 0.9079 & 0.8452 & $-0.0627$ \\
\bottomrule
\end{tabular}
\end{table}

Three of the four classes transfer with little or no loss, and
\texttt{coastal\_mansion} in fact improves. Essentially the entire 6.3-point
macro-$F_1$ drop is attributable to \texttt{dense\_residential}, which falls by
19 points. On the tuning split the model appears to have largely solved that
class at $F_1 = 0.917$; on unseen data it has not, at $F_1 = 0.727$. The
selected checkpoint is thus the epoch at which the model best fitted the
particular density boundary present in those 420 tuning images, and that
boundary does not generalise.

This is a concrete illustration of why Section~\ref{sec:splits} insists on a
sealed held-out set. Had the tuning metrics been reported as final --- a
natural mistake, since they are computed on data the model never received
gradients from --- we would have overstated this model's performance by more
than six points and, more seriously, drawn the wrong conclusion about which
class remains unsolved.

The training curves (Figure~\ref{fig:customcnn-curves}) reveal two distinct
phases. Through roughly the first fifteen epochs the model exhibits pronounced
validation instability, with tuning accuracy oscillating between 0.25 and 0.90
and validation loss spiking above 4.0, while training accuracy rises smoothly
throughout. We attribute this asymmetry to the batch-normalisation moving
averages, which are estimated from batch statistics during training but applied
as running averages at validation time, and which converge slowly on a
420-image tuning split. After the first learning-rate reduction at epoch 16 the
oscillation collapses and validation loss flattens, giving a clear visual record
of the scheduler earning its place in the configuration.

\begin{figure}[htbp]
\centering
\includegraphics[width=\linewidth]{docs/figures/custom_cnn_curves.png}
\caption{Custom CNN training curves. The dashed line marks epoch 24, the
checkpoint selected on tuning accuracy. Early validation instability is
consistent with batch-normalisation statistics converging on a small tuning
split; it subsides once the learning-rate scheduler first fires at epoch 16.
The persistent gap between the training and validation curves after epoch 20
indicates residual overfitting.}
\label{fig:customcnn-curves}
\end{figure}

The second phase exposes the model's central limitation. Training accuracy
continues to climb toward 0.99 while validation accuracy plateaus near 0.87 and
validation loss ceases to improve, leaving a gap of roughly twelve points. The
L2 term reduces this overfitting but does not eliminate it, which is
unsurprising given that 25.7M of the 26.5M parameters reside in the first dense
layer alone --- 97\% of the model's capacity is concentrated in a single layer
fitted to 2,378 images.

\subsubsection{Discussion}

Our result does not reproduce the ordering reported in \cite{Abdelmonem2026},
where the custom architecture is presented as competitive with or superior to
pretrained networks. On our partition it is decisively outperformed by all four
transfer-learning approaches, including MobileNetV2, which reaches 0.9576
macro-$F_1$ with 2.26M parameters --- roughly one twelfth of the custom model's
capacity.

We attribute the discrepancy to the granularity of the classification task
rather than to any defect in the architecture. The ten NWPU-RESISC45 categories
used in the source study --- harbour, chaparral, tennis court, industrial area,
parking lot, forest, beach, overpass, airplane, and baseball diamond --- are
semantically and visually disjoint, and are separable on coarse cues such as
dominant colour, the presence of water, or large geometric structures. Features
sufficient for those distinctions are learnable from 700 images per class. Our
four categories are all residential or institutional building scenes
distinguished by building density and spatial arrangement, and the
discriminative features are correspondingly subtler. ImageNet pretraining
supplies exactly the mid-level texture and structure detectors that such a task
demands and that 2,378 images are insufficient to induce from scratch.

It is also worth noting that the source paper's own comparison is more equivocal
than its abstract suggests: on NWPU-RESISC45 (its Table~7) VGG16 attains 0.949
and Xception 0.948 test accuracy against the proposed model's 0.9428, so the
architecture's advantage there rests on its efficiency and its markedly stronger
performance on the smaller UC~Merced dataset rather than on raw accuracy. Read
in that light, our finding is less a contradiction of \cite{Abdelmonem2026} than
a delineation of where its conclusions hold: a well-designed from-scratch CNN is
competitive on coarse-grained scene classification, but pretraining remains
decisive when categories are fine-grained and data is limited. Establishing that
boundary is precisely the contribution this baseline makes to our comparison,
and it is developed further in Section~\ref{sec:comparative}.
\subsection{ResNet-50 (transfer learning)}
% 23.6M total / 8.2K trainable — frozen ImageNet backbone + new head.
ResNet's key idea is the residual (skip) connection: each block learns an
adjustment to its input rather than a full transformation, which lets
gradients flow through many layers and made genuinely deep networks trainable
for the first time \cite{He2016}. ResNet-50, pretrained on ImageNet, is one of
the most widely used backbones in computer vision, with transfer behaviour
that is well understood making it one of the natural strong defaults among our
five approaches.

We use the backbone as a frozen feature extractor and train only a small new
classification head. The whole pipeline is packaged inside a single Keras
model (Table~\ref{tab:resnet50-config}): the saved checkpoint accepts raw
$[0, 255]$ RGB images and handles augmentation and preprocessing itself.

\begin{table}[H]
\centering
\caption{ResNet-50 model configuration.}
\label{tab:resnet50-config}
\small
\begin{tabular}{ll}
\toprule
Stage & Details \\
\midrule
Input & $256 \times 256 \times 3$, raw $[0,255]$ RGB \\
Augmentation & shared team policy, active during training only \\
Preprocessing & ResNet \texttt{preprocess\_input} (ImageNet channel normalisation) \\
Backbone & ResNet-50, ImageNet weights, frozen, global average pooling \\
Head & dropout 0.3 $\rightarrow$ dense softmax (4 units) \\
Parameters & 23{,}595{,}908 total; 8{,}196 trainable \\
Optimisation & Adam, learning rate $10^{-3}$, categorical cross-entropy \\
\bottomrule
\end{tabular}
\end{table}
Three design decisions are worth spelling out. First, the backbone is frozen rather than fine-tuned: with 2,378 training images against 23.6M parameters, full fine-tuning risks overfitting, and freezing matches what the other transfer models do, keeping the comparison clean. Only the 8,196 parameter head is fitted to our data. Second, augmentation and preprocessing live inside the model graph rather than in the data
pipeline, so the deployed checkpoint cannot be fed wrongly normalised inputs
--- what the model sees at inference is exactly what it saw in training.
Third, the two head hyper-parameters (dropout rate and learning rate) are
exposed as controlled experiment factors; the chosen values (0.3 and
$10^{-3}$) were selected on the internal tuning split and never against the
held-out set.

\subsubsection{Results}

Training used the shared protocol --- batch size 16, up to 32 epochs, early
stopping with patience 20 --- and selected the best checkpoint at epoch 25,
after 2,043 seconds. On the held-out set the model achieves \textbf{97.5\%
accuracy} and \textbf{0.9750 macro-$F_1$}, the second strongest of our five
approaches. Per-class $F_1$ ranges narrowly from 0.964
(\texttt{dense\_residential}) to 0.995 (\texttt{coastal\_mansion}), so no class
is left behind. The full leaderboard is given in Section~\ref{sec:results}.

\begin{figure}[htbp]
\centering
\includegraphics[width=\linewidth]{docs/figures/resnet50_curves.png}
\caption{ResNet-50 training curves. The dashed line marks epoch 25, the
selected checkpoint. Training and validation curves track each other closely
throughout, and validation accuracy frequently exceeds training accuracy
because augmentation is active only on the training stream.}
\label{fig:resnet50-curves}
\end{figure}

The curves in Figure~\ref{fig:resnet50-curves} differ markedly from those of
the from-scratch model. Validation accuracy exceeds 0.91 after a single epoch
and settles above 0.95 within five, reflecting that the frozen ImageNet
features are already discriminative for this task and only the 8,196-parameter
head has to be fitted. There is no widening gap between the training and
validation curves; indeed validation accuracy frequently sits \emph{above}
training accuracy, which is expected here because augmentation is applied to
the training stream only, so the model is scored on clean images at validation
time but on flipped, rotated, and jittered ones during training. With so few
trainable parameters the model has almost no capacity to memorise the training
set, and the residual oscillation around epochs 13--16 reflects the small size
of the 420-image tuning split rather than instability in the model.

\begin{figure}[htbp]
\centering
\begin{minipage}[t]{0.49\textwidth}
  \centering
  \includegraphics[width=\linewidth]{docs/figures/resnet50_tuning_confusion_matrix.png}
  \\[0.2em] {\small (a) Tuning split (420 images)}
\end{minipage}\hfill
\begin{minipage}[t]{0.49\textwidth}
  \centering
  \includegraphics[width=\linewidth]{docs/figures/resnet50_confusion_matrix.png}
  \\[0.2em] {\small (b) Held-out set (400 images)}
\end{minipage}
\caption{ResNet-50 confusion matrices. Unlike the from-scratch model, the two
panels are closely similar: the frozen backbone leaves little room for the
checkpoint to overfit the tuning split.}
\label{fig:resnet50-cm}
\end{figure}

\paragraph{Checkpoint-selection bias is negligible.}
Repeating the tuning-versus-held-out comparison made for the custom CNN
(Table~\ref{tab:resnet50-bias}) produces a very different picture. Every class
transfers within $\pm 0.02$ $F_1$, and the macro-$F_1$ falls by only 0.006 ---
roughly one tenth of the 0.063 drop seen in the from-scratch model.

\begin{table}[H]
\centering
\caption{ResNet-50 per-class $F_1$ on the tuning split versus the held-out set.}
\label{tab:resnet50-bias}
\small
\begin{tabular}{lccr}
\toprule
Class & Tuning & Held-out & $\Delta$ \\
\midrule
\texttt{coastal\_mansion}    & 0.977 & 0.995 & $+0.018$ \\
\texttt{sparse\_residential} & 0.976 & 0.966 & $-0.010$ \\
\texttt{nursing\_home}       & 0.990 & 0.975 & $-0.015$ \\
\texttt{dense\_residential}  & 0.981 & 0.964 & $-0.017$ \\
\midrule
Macro-$F_1$ & 0.9810 & 0.9750 & $-0.0059$ \\
\bottomrule
\end{tabular}
\end{table}

The contrast is instructive, and it is explained by trainable capacity rather
than by architecture family. Selecting a checkpoint on tuning accuracy can only
inflate the reported score to the extent that the model is able to fit the
idiosyncrasies of those 420 images. ResNet-50 adapts 8,196 parameters and shows
a 0.6-point optimism; the custom CNN adapts 26,511,444 and shows a 6.3-point
optimism concentrated almost entirely in one class. Freezing the backbone
therefore buys not only accuracy but also \emph{trustworthiness of the
validation signal} --- a benefit that a leaderboard alone does not reveal. This
relationship is examined across all five models in
Section~\ref{sec:comparative}.

\subsection{EfficientNet-B0 (transfer learning)}
\todo{}
\subsection{MobileNetV2 (transfer learning)}
% 2.3M total / 5.1K trainable — efficiency-oriented baseline.
\todo{}
\subsection{Vision Transformer (ViT)}
\todo{Incomplete}
% 85.8M total / 3.1K trainable; input resize fix; attention-based contrast
% to the convolutional models.
The Vision Transformer \cite{Dosovitskiy2021} applies a pure Transformer
architecture to images: the input is split into fixed $16\times16$ patches,
each patch is linearly embedded and given a learned position embedding, and
the resulting sequence is processed by a stack of self-attention layers. Unlike a CNN, whose
convolutions can only relate nearby pixels layer by layer, self-attention
lets every patch attend to every other patch from the first layer onward,
capturing global layout cues directly. This makes ViT a structurally
different approach from our four convolutional models, which is exactly the
contrast our comparison needs: aerial scene categories such as dense versus
sparse residential are distinguished less by any single local texture than
by the overall spatial arrangement of buildings, roads, and vegetation.

We use the ViT-B/16 backbone pretrained on ImageNet as a frozen feature
extractor and train only a small classification head on the CLS token
embedding. As with the other approaches, the whole pipeline is packaged
inside a single Keras model (Table~\ref{tab:vit-config}): the saved
checkpoint accepts raw $[0,255]$ RGB images and applies augmentation and
preprocessing itself.

\begin{table}[H]
\centering
\caption{ViT-B/16 model configuration.}
\label{tab:vit-config}
\small
\begin{tabular}{ll}
\toprule
Stage & Details \\
\midrule
Input & $256\times256\times3$, raw $[0,255]$ RGB \\
Augmentation & shared team policy, active during training only \\
Resize & $256\times256 \rightarrow 224\times224$ (bilinear), see below \\
Preprocessing & scale to $[-1,1]$ (ViT convention, not ImageNet channel stats) \\
Backbone & ViT-B/16, ImageNet weights, frozen; 12 layers, 768 hidden dim, 12 heads \\
Head & CLS token $\rightarrow$ dropout 0.3 $\rightarrow$ dense softmax (4 units) \\
Parameters & 85{,}801{,}732 total; 3{,}076 trainable \\
Optimisation & Adam, learning rate $10^{-3}$, categorical cross-entropy \\
\bottomrule
\end{tabular}
\end{table}

The model resizes its input from the shared $256\times256$ to $224\times224$ immediately before the backbone. 
This is not optional: ViT's learned position embeddings form
a fixed table with one entry per patch position ($14\times14$ patches
$+$ CLS $=$ 197 tokens for a $224\times224$ input), so a $256\times256$
input would produce 257 tokens with no pretrained embedding to match. CNNs
tolerate off-resolution input silently because convolution weights are
shared across positions; ViT makes the mismatch explicit. Notably, the CNN
backbones used by the other approaches were also pretrained at
$224\times224$, so our ViT is arguably the only model evaluated at its
native pretraining resolution. The resize is placed after the shared
augmentation, so regularisation is applied identically to all five models.
Second, the backbone is frozen: ViTs lack the inductive biases of
convolutions and are known to need large datasets when fine-tuned, so with
only 2{,}798 usable training images we fit just the 3{,}076-parameter head
--- by far the smallest trainable footprint of our five approaches, which
also makes overfitting nearly impossible. Third, preprocessing scales
pixels to $[-1,1]$ rather than applying ImageNet channel normalisation,
matching how this checkpoint was pretrained.

\paragraph{Pretraining-source ablation.}
The choice of pretrained weights was validated experimentally rather than
assumed. The selected checkpoint was pretrained on ImageNet-21k and then
finetuned on ImageNet-1k; we compared it against the raw ImageNet-21k
weights (never finetuned to the 1k label space) under an otherwise identical
protocol --- same seed, split, augmentation, and training flags. Our
hypothesis was that the more generic 21k features would transfer better to
aerial imagery, which is a clear domain shift from natural photos. To keep
the held-out set untouched by model selection, the comparison used the
internal tuning split only (Table~\ref{tab:vit-ablation}).

\begin{table}[H]
\centering
\caption{ViT pretraining-source ablation on the internal tuning split.}
\label{tab:vit-ablation}
\small
\begin{tabular}{lccc}
\toprule
Variant & Tuning accuracy & Tuning macro-$F_1$ & Best epoch \\
\midrule
ImageNet-1k finetuned (selected) & 0.9548 & 0.9549 & 10 \\
ImageNet-21k raw                 & 0.9405 & 0.9408 & 31 \\
\bottomrule
\end{tabular}
\end{table}

The hypothesis was cleanly falsified: the 1k-finetuned checkpoint scored
1.4 points higher macro-$F_1$ and converged three times faster (best epoch
10 vs.\ 31, with the 21k variant still improving when the 32-epoch budget
ended). The explanation lies in how the features are used. Supervised
finetuning on the 1k label space shapes the CLS embedding to be linearly
separable --- precisely what a frozen-backbone setup with a softmax head
requires. The raw 21k features are not poorer, but more entangled: a
3{,}076-parameter head can only draw four hyperplanes through the
768-dimensional embedding space and cannot re-mix features, so it benefits
from geometry that finetuning has already organised. The slow crawl of the
21k run is the visible signature of that harder optimisation problem. With
full fine-tuning the outcome could well reverse; in the frozen regime, the
linear-separability effect dominated the domain-shift effect. The 1k
variant was therefore retained, and only it was evaluated on the held-out
set.

\subsubsection{Results}

Training used the shared protocol --- batch size 16, up to 32 epochs, early
stopping with patience 20 --- and selected the best checkpoint at epoch 10,
the earliest of our five approaches; training stopped at epoch 30 after
9{,}378 seconds on CPU. On the held-out set the model achieves
\textbf{96.5\% accuracy} and \textbf{0.9652 macro-$F_1$}, third of the five
approaches and within 1.25 points of the strongest. Per-class $F_1$ ranges
from 0.951 (\texttt{sparse\_residential} and \texttt{dense\_residential})
to 0.980 (\texttt{nursing\_home} and \texttt{coastal\_mansion}, the latter
with perfect precision --- no image of any other class was ever predicted
as a coastal mansion). The full leaderboard is given in
Section~\ref{sec:results}.

\begin{figure}[htbp]
\centering
\includegraphics[width=\linewidth]{docs/figures/ViT_curves.png}
\caption{ViT training curves. The dashed line marks epoch 10, the selected
checkpoint. Validation accuracy exceeds 0.87 after a single epoch and
plateaus within ten: the frozen ViT features are immediately usable and the
3{,}076-parameter head has little left to learn thereafter.}
\label{fig:vit-curves}
\end{figure}

The curves in Figure~\ref{fig:vit-curves} show the fastest convergence of
our five approaches, consistent with the ablation's interpretation: the
1k-finetuned CLS embedding is already organised for linear classification,
so the head locks on within a few epochs and the remaining twenty epochs
add only noise around the plateau. Training and validation accuracy track
each other closely throughout, with no widening gap.

\begin{figure}[htbp]
\centering
\begin{minipage}[t]{0.49\textwidth}
  \centering
  \includegraphics[width=\linewidth]{docs/figures/ViT_tuning_confusion_matrix.png}
  \\[0.2em] {\small (a) Tuning split (420 images)}
\end{minipage}\hfill
\begin{minipage}[t]{0.49\textwidth}
  \centering
  \includegraphics[width=\linewidth]{docs/figures/ViT_confusion_matrix.png}
  \\[0.2em] {\small (b) Held-out set (400 images)}
\end{minipage}
\caption{ViT confusion matrices. Errors concentrate on the
dense/sparse-residential boundary in both panels; on the held-out set the
single largest error mode is 4 \texttt{sparse\_residential} images predicted
as \texttt{dense\_residential}.}
\label{fig:vit-cm}
\end{figure}

The confusion matrices (Figure~\ref{fig:vit-cm}) repeat the pattern seen
across all five models: the dense-versus-sparse residential boundary --- a
density judgement rather than a structural one --- absorbs most of the
mistakes, while \texttt{coastal\_mansion} and \texttt{nursing\_home} are
essentially never confused with each other.

\paragraph{Checkpoint-selection bias is absent.}
Repeating the tuning-versus-held-out comparison made for the previous
models (Table~\ref{tab:vit-bias}), the ViT sits at the far end of the
capacity spectrum: with only 3{,}076 trainable parameters, macro-$F_1$ is
actually \emph{higher} on the held-out set than on the tuning split used to
select the checkpoint ($+0.010$). The selection optimism visible in the
custom CNN (6.3 points) and faintly in ResNet-50 (0.6 points) has here
fallen below zero --- the model cannot meaningfully fit the idiosyncrasies
of 420 tuning images, so the tuning signal transfers essentially unbiased.
This completes the capacity-versus-trustworthiness relationship examined in
Section~\ref{sec:comparative}.

\begin{table}[H]
\centering
\caption{ViT per-class $F_1$ on the tuning split versus the held-out set.}
\label{tab:vit-bias}
\small
\begin{tabular}{lccr}
\toprule
Class & Tuning & Held-out & $\Delta$ \\
\midrule
\texttt{coastal\_mansion}    & 0.950 & 0.980 & $+0.030$ \\
\texttt{nursing\_home}       & 0.971 & 0.980 & $+0.009$ \\
\texttt{sparse\_residential} & 0.947 & 0.951 & $+0.004$ \\
\texttt{dense\_residential}  & 0.952 & 0.951 & $-0.001$ \\
\midrule
Macro-$F_1$ & 0.9549 & 0.9652 & $+0.0103$ \\
\bottomrule
\end{tabular}
\end{table}

\section{Model Training and Tuning}
% The shared protocol that makes the comparison FAIR: identical seed (42),
% batch size 16, image size 256, max 32 epochs, early stopping patience 20,
% best-epoch checkpointing on tuning accuracy. Per-model hyper-parameters
% (dropout 0.3, lr 1e-3 where applicable) and how they were chosen.
% Include training/validation curves (results/<name>_curves.png) and comment
% on convergence: e.g. ViT best epoch 10 vs MobileNetV2 best epoch 30.
% Hardware note: mixed GPU (WSL) and CPU training; wall-clock times.
\todo{Training protocol + curves}

\section{Inference and Performance Evaluation}
\subsection{Evaluation Protocol}
% Fixed-schema evaluate(): accuracy, per-class and macro precision/recall/F1,
% confusion matrix, on the 400 held-out images (100/class).
\subsection{Results}
\label{sec:results}
% Held-out (val 23) leaderboard — numbers below are from results/*.json:
\begin{table}[H]
\centering
\caption{Held-out performance on \texttt{val 23} (400 images, 100 per class).}
\label{tab:leaderboard}
\begin{tabular}{lccccrr}
\toprule
Model & Accuracy & Macro P & Macro R & Macro F1 & Params & Best Epoch \\
\midrule
ResNet-50        & \textbf{0.9750} & 0.9757 & 0.9750 & \textbf{0.9750} & 23.6M & fillup \\
ViT              & 0.9650 & 0.9659 & 0.9650 & 0.9652 & 85.8M & fillup \\
MobileNetV2      & 0.9575 & 0.9592 & 0.9575 & 0.9576 & 2.3M  & fillup  \\
Custom CNN       & 0.8450 & 0.8772 & 0.8450 & 0.8452 & 26.5M & fillup  \\
EfficientNet-B0  & \todo{} & \todo{} & \todo{} & \todo{} & \todo{} & \todo{} \\
\bottomrule
\end{tabular}
\end{table}
% Confusion matrices (results/<name>_confusion_matrix.png) side by side.
% Per-class analysis: the custom CNN collapses on dense_residential
% (recall 0.64, mostly confused with sparse_residential) while transfer
% models resolve it — a strong "why pretraining matters" narrative.
\todo{Figures + per-class discussion}

\section{Strengths, Weaknesses, and Comparative Insights}
\label{sec:comparative}
% Required by the spec: discuss each approach EXTENSIVELY.
% Suggested axes: accuracy vs parameter count vs training cost vs inference
% latency; data efficiency of transfer learning; ViT's cost/benefit at this
% dataset scale; MobileNetV2 as the best accuracy-per-parameter;
% error overlap between models (do they fail on the same images?).
\todo{Comparative analysis table + discussion}

\section{Deployment via Containerisation}
% Docker image: python:3.12-slim, Streamlit UI (frontend/web.py) on port 8501,
% healthcheck, best-model checkpoint baked in. Upload an aerial image ->
% predicted label + confidence (softmax). Include architecture diagram,
% build/run commands, and UI screenshot.
% NOTE: deploy the BEST model — EfficientNet-B0 (0.9775 held-out macro-F1).
% The dockerfile currently copies custom_cnn_best.keras, so both the COPY line,
% .dockerignore allow-list, and frontend/web.py MODEL_PATH must be switched to
% efficientnet_b0_best.keras before the final build.
\todo{Deployment description + screenshots}

\section{Results, Insights, and Discussion}
% Synthesis: what the benchmark says about fine-grained aerial classification;
% final model recommendation and justification (accuracy vs efficiency);
% limitations (single seed, no cross-validation, 400-image test set ->
% ~±2% confidence interval on accuracy).
\todo{Key insights}

%=====================================================================
% 3. INDIVIDUAL CONTRIBUTIONS  (1–2 pages)
%=====================================================================
\section{Individual Contributions}
\label{sec:contributions}
% One short paragraph per member: model owned, plus shared-infrastructure /
% EDA / deployment / report work. Keep it concrete (branches: model/customcnn,
% model/ResNet-50, model/effNet-b0, mobilenetv2, model/vit, feature/frontend).
\renewcommand{\arraystretch}{1.35}
\begin{tabular}{@{}l@{\hspace{2.5em}}l@{}}
\toprule
\textbf{Member} & \textbf{Contribution} \\
\midrule
Pang Jia De         & Custom CNN --- from-scratch baseline \\
Daryl Chong Han Kin & ResNet-50 --- transfer learning \\
Tan Jian Xin        & MobileNetV2 --- transfer learning \\
Kaiser               & Vision Transformer (ViT-B/16) --- transfer learning \\
Jun Kiat            & EfficientNet-B0 --- transfer learning \\
\bottomrule
\end{tabular}
\renewcommand{\arraystretch}{1}
\todo{need 1-2 page rn too little}
%=====================================================================
% 4. REFLECTION  (2–3 pages)
%=====================================================================
\section{Reflection}
% What you learnt and how you applied it: e.g. leakage-safe evaluation design,
% controlled experiments (fixed batch size/seed across members), transfer
% learning in low-data regimes, team engineering (frozen shared/ contract,
% branch workflow, result schemas), debugging (ViT input resize fix,
% TF Python-version constraints, WSL GPU setup).


%=====================================================================
% 5. FUTURE WORK  (1–2 pages)
%=====================================================================
\section{Future Work}
% e.g. fine-tuning (unfreezing) backbones with discriminative learning rates,
% ensembling the top models, test-time augmentation, Grad-CAM interpretability,
% cross-validation over seeds, REST inference API (FastAPI) alongside the UI,
% expanding to more of the 38-category benchmark.
\todo{Future directions}

\bibliography{mybib}

\appendix
\section{Supplementary Material}
% Full hyper-parameter tables, extra confusion matrices, additional EDA plots,
% repository structure, reproduction instructions (train.py command line).
\todo{Appendix}

\end{document}

